\chapter{Python}\label{ch:iv:python}

A helper that appends each fill to a default list argument works in the unit test and doubles every position by the end of
the first live day. The candidate asked about it in an interview has twenty seconds to say why the list was shared, and another
minute to say how he would have found it without being told. Python interviews for researchers and developers test the data
model (names, objects, mutability), iteration, the interpreter's concurrency model, and the numerical stack's views, copies and
vectorisation; the platform side is One Quant Book~15, chapter~8. Every snippet here runs in the chapter's tests with the pinned
interpreter (CPython~3.10), and where an answer depends on CPython rather than on the language, the solution says so.

\section{The data model: names, objects and mutability}

A Python variable is a name bound to an object; assignment rebinds a name and never copies. Mutable objects (lists, dicts,
sets, most user classes) can change in place, and every name bound to them sees the change. Default argument values are
evaluated once, when the \texttt{def} statement runs, so a mutable default is shared by every call. Equality (\texttt{==}) compares
values through \texttt{\_\_eq\_\_}; identity (\texttt{is}) compares objects, and whether two equal immutable values are the same
object is an implementation detail. A type used as a dictionary key needs \texttt{\_\_hash\_\_} consistent with \texttt{\_\_eq\_\_} and
must not change while it is a key, which is what a frozen dataclass provides.

\omcode{code/interviews/24-python/python/snippets/iv_snip_mutable_default.py}{1}{8}{A mutable default
argument.}\label{lst:iv:python:default}

\section{Iterators, generators and context managers}

An iterator produces values on demand and is exhausted after one pass; a generator is an iterator written as a function with
\texttt{yield}, or as a generator expression. Pipelines of generators process a file of any size in constant memory, because
each stage pulls one record at a time from the one before. A context manager pairs set-up with guaranteed tear-down
(\texttt{with}), which is how files, locks and temporary changes of state are made safe against exceptions.

\omcode{code/interviews/24-python/python/iv_py.py}{8}{25}{A generator pipeline over a file of fills: read, filter,
aggregate, one line in memory at a time.}\label{lst:iv:python:pipeline}

Closures capture variables, not values: a function created in a loop sees the loop variable's final value when called later
(\cref{iq:iv:python:2}). Mutating a dictionary while iterating over it raises an error; build a new one instead.

\section{The interpreter lock, threads, processes and async}

CPython's global interpreter lock (One Quant Book~15, chapter~8) lets one thread execute Python bytecode at a time. Threads
therefore speed up work that waits (network, disk) or that runs in C code releasing the lock (most of numpy's heavy
operations, compression, many I/O libraries), and do not speed up pure-Python computation, for which processes are the answer,
at the cost of serialising data between them. \texttt{asyncio} runs many waiting tasks on one thread with explicit yield points,
suited to many network connections (a crypto market-data client with hundreds of websocket streams). The lock does not make
compound operations atomic: \texttt{x += 1} on a shared variable is several bytecodes, and threads can interleave between them.

\section{The numerical stack}

Vectorisation replaces a Python loop by one call into compiled code over an array, typically tens to hundreds of times faster
for arithmetic (One Quant Book~15, chapter~8, on interpreter overhead and array programming). Three rules avoid the classic bugs.
Basic slicing of a numpy array returns a view that shares memory with the original; fancy indexing (a list or boolean mask)
returns a copy. In pandas, assign through \texttt{.loc} in one step: chained indexing such as \texttt{df[mask]['col'] = v} may assign
to a temporary copy and leave the frame unchanged. And floating-point sums depend on order: \texttt{math.fsum} or pairwise
summation (numpy's) keeps error small where a naive running sum does not (One Quant Book~4, chapter~25).

\omcode{code/interviews/24-python/python/iv_py.py}{43}{46}{A rolling mean by cumulative sums: linear time and no Python
loop.}\label{lst:iv:python:rolling}

\begin{method}[Answering a Python output question]\label{met:iv:python:output}
\begin{enumerate}
\item Track names and the objects they are bound to; ask which objects are mutable and shared.
\item Note when each expression is evaluated: defaults at definition, closures at call, generators on demand.
\item Separate language guarantees from CPython details (small-integer caching, dictionary internals).
\item For numpy and pandas, ask for every indexing operation whether it returns a view or a copy.
\end{enumerate}
\end{method}

\section{Question bank}

\begin{interviewq}[$\star$ \iqroles{developer, researcher} \iqfirm{systematic fund}]\label{iq:iv:python:1}
What does \cref{lst:iv:python:default} print, and how do you fix the function?
\end{interviewq}

\omcode{code/interviews/24-python/python/snippets/iv_snip_late_binding.py}{1}{4}{Pricers made in a
loop.}\label{lst:iv:python:late}

\begin{interviewq}[$\star$ \iqroles{developer, researcher} \iqfirm{any}]\label{iq:iv:python:2}
What do the two lines of \cref{lst:iv:python:late} print, and why does the second fix the first?
\end{interviewq}

\omcode{code/interviews/24-python/python/snippets/iv_snip_identity.py}{1}{6}{Equality and
identity.}\label{lst:iv:python:identity}

\begin{interviewq}[$\star$ \iqroles{developer} \iqfirm{any}]\label{iq:iv:python:3}
What does \cref{lst:iv:python:identity} print under CPython? Which parts of that answer does the language guarantee?
\end{interviewq}

\omcode{code/interviews/24-python/python/snippets/iv_snip_generator_twice.py}{1}{3}{A generator used
twice.}\label{lst:iv:python:gen}

\begin{interviewq}[$\star$ \iqroles{researcher, mle} \iqfirm{systematic fund}]\label{iq:iv:python:4}
What does \cref{lst:iv:python:gen} print?
\end{interviewq}

\omcode{code/interviews/24-python/python/snippets/iv_snip_dict_mutation.py}{1}{9}{Removing flat
positions.}\label{lst:iv:python:dict}

\begin{interviewq}[$\star\star$ \iqroles{developer, researcher} \iqfirm{multi-manager fund}]\label{iq:iv:python:5}
What happens in the loop of \cref{lst:iv:python:dict}, and why is the last line the right way to do it?
\end{interviewq}

\begin{interviewq}[$\star\star$ \iqroles{developer, researcher} \iqfirm{systematic fund}]\label{iq:iv:python:6}
Compute notional by symbol over the large fills in a 50-gigabyte CSV file on a machine with 16 gigabytes of memory, in plain
Python.
\end{interviewq}

\begin{interviewq}[$\star\star$ \iqroles{developer} \iqfirm{bank}]\label{iq:iv:python:7}
Write a context manager that changes a risk limit for the duration of a block and restores it even if the block raises.
\end{interviewq}

\omcode{code/interviews/24-python/python/snippets/iv_snip_views.py}{1}{9}{Views and copies.}\label{lst:iv:python:views}

\begin{interviewq}[$\star\star$ \iqroles{researcher, mle} \iqfirm{systematic fund}]\label{iq:iv:python:8}
What does \cref{lst:iv:python:views} print?
\end{interviewq}

\omcode{code/interviews/24-python/python/snippets/iv_snip_float_sum.py}{1}{5}{Summing ten ticks of
0.1.}\label{lst:iv:python:fsum}

\begin{interviewq}[$\star\star$ \iqroles{researcher, developer} \iqfirm{market maker}]\label{iq:iv:python:9}
What does \cref{lst:iv:python:fsum} print? When does the difference matter on a desk?
\end{interviewq}

\begin{interviewq}[$\star\star\star$ \iqroles{developer, mle} \iqfirm{crypto firm}]\label{iq:iv:python:10}
A service subscribes to 300 websocket market-data streams and computes a statistic per message. When do threads help, when
do processes, and when \texttt{asyncio}? What does the interpreter lock not protect you from?
\end{interviewq}

\begin{interviewq}[$\star\star\star$ \iqroles{researcher, mle} \iqfirm{systematic fund}]\label{iq:iv:python:11}
Vectorise a rolling mean of window $k$ over a price array without a Python loop, and say how you check it.
\end{interviewq}

\omcode{code/interviews/24-python/python/snippets/iv_snip_chained.py}{1}{12}{Flagging short positions in
pandas.}\label{lst:iv:python:chained}

\begin{interviewq}[$\star\star\star$ \iqroles{researcher, developer} \iqfirm{multi-manager fund}]\label{iq:iv:python:12}
What does \cref{lst:iv:python:chained} print, and why? How does copy-on-write, announced as the default for pandas 3.0, change the answer?
\end{interviewq}

\begin{interviewq}[$\star\star\star$ \iqroles{developer} \iqfirm{proprietary firm}]\label{iq:iv:python:13}
Write a value type for an order key (venue, client order identifier) that can be used in a dictionary. What goes wrong if a
mutable object is used as a key and then changed?
\end{interviewq}

\begin{omsources}
\item The Python Language Reference and Library Reference, version 3.10 (the series' pinned interpreter); the pandas 2.3 and
numpy 2.2 user guides on indexing, views and copies, and copy-on-write.
\item One Quant Book~15, chapter~8 (Python at scale); One Quant Book~4, chapter~25 (floating point).
\end{omsources}
